% BACK fresh Korean production translation.
% Initial translation authority: ED0005 / NOETH-DE-AUTH-v038-20260805.
% Current scope-rebase authority: ED0011 / NOETH-DE-AUTH-v047-20260808.
% Complete translated scope: current German lines 23984--24146.

\clearpage
\begin{center}
{\Large\bfseries 참고문헌}
\end{center}

\begin{enumerate}
\setlength{\itemsep}{0.35em}
\item 삼원 사차형식의 형식계 구성에 관하여.\\
Sitz. Ber. d. Physikal.-mediz. Sozietät in Erlangen 39 (1907), 176--179쪽
\item 삼원 사차형식의 형식계 구성에 관하여.\\
J. Reine Angew. Math. 134 (1908), 23--90쪽 및 표 1개
\item \(n\)개 변수 형식의 불변식론에 관하여.\\
J. Ber. d. DMV 19 (1910), 101--104쪽
\item \(n\)개 변수 형식의 불변식론에 관하여.\\
J. Reine Angew. Math. 139 (1911), 118--154쪽
\item 유리함수체.\\
J. Ber. d. DMV 22 (1913), 316--319쪽
\item 체와 유리함수계.\\
Math. Ann. 76 (1915), 161--196쪽
\item 유한군 불변식의 유한성 정리.\\
Math. Ann. 77 (1916), 89--92쪽
\item 임의로 많은 기본형식들의 계의 불변식에 대한 다항식적 표현.\\
Math. Ann. 77 (1916), 93--102쪽. (16번도 참조)
\item 정수적 초월수로 이루어진 가장 일반적인 영역.\\
Math. Ann. 77 (1916), 103--128쪽. (16번도 참조)
\item 동형사상의 함수방정식.\\
Math. Ann. 77 (1916), 536--545쪽
\item 지정된 군을 갖는 방정식.\\
Math. Ann. 78 (1918), 221--229쪽. (16번도 참조)
\item 임의의 미분식의 불변식.\\
Nachr. v. d. Ges. d. Wiss. zu Göttingen 1918, 37--44쪽
\item 불변 변분문제.\\
Nachr. v. d. Ges. d. Wiss. zu Göttingen 1918, 235--257쪽
\item 한 변수 대수함수의 산술이론과 다른 이론들 및 수체론과의 관계.\\
J. Ber. d. DMV 28 (1919), 182--203쪽
\item 이원형식의 정수 계수 불변식계의 유한성.\\
Nachr. v. d. Ges. d. Wiss. zu Göttingen 1919, 138--156쪽
\item 형식론의 급수전개에 관하여.\\
Math. Ann. 81 (1920), 25--30쪽
\item W. 슈마이들러와 공동: 비가환 영역의 가군, 특히 미분식 및 차분식으로 이루어진 가군.\\
Math. Z. 8 (1920), 1--35쪽
\item 전쟁 중 전사한 K. 헨첼트의 소거이론 논문에 관하여.\\
J. Ber. d. DMV 30 (1921), 48쪽
\item 환영역의 아이디얼 이론.\\
Math. Ann. 83 (1921), 24--66쪽
\item 절대기약성의 한 대수적 판정기준.\\
Math. Ann. 85 (1922), 26--33쪽
\item 형식적 변분법과 미분불변식.\\
Encyklopädie d. math. Wiss. III, 3 (1922), 68--71쪽(R. Weitzenböck, Differentialinvarianten 안에 수록)
\item K. 헨첼트 \(\dagger\)의 원고 편집: 다항식 아이디얼과 종결식 이론.\\
Math. Ann. 88 (1923), 53--79쪽
\item 대수불변식과 미분불변식.\\
J. Ber. d. DMV 32 (1923), 177--184쪽
\item 소거이론과 일반 아이디얼 이론.\\
Math. Ann. 90 (1923), 229--261쪽
\item 소거이론과 아이디얼 이론.\\
J. Ber. d. DMV 33 (1924), 116--120쪽
\item 대수적 수체에서의 아이디얼 이론의 추상적 구성.\\
J. Ber. d. DMV 33 (1924), 102쪽
\item 아이디얼 이론의 힐베르트 수.\\
J. Ber. d. DMV 34 (1925), 101쪽
\item 군 지표와 아이디얼 이론.\\
J. Ber. d. DMV 34 (1925), 144쪽
\item 표수 \(p\)인 유한 선형군의 불변식에 관한 유한성 정리.\\
Nachr. v. d. Ges. d. Wiss. zu Göttingen 1926, 28--35쪽
\item 대수적 수체와 함수체에서 아이디얼 이론의 추상적 구성.\\
Math. Ann. 96 (1927), 26--61쪽
\item 대수적 수체 또는 함수체의 오더에 대한 판별식 정리.\\
J. Reine Angew. Math. 157 (1927), 82--104쪽
\item R. 브라우어와 공동: 기약표현의 최소 분해체에 관하여.\\
Sitz. Ber. d. Preuß. Akad. d. Wiss. 1927, 221--228쪽
\item 산술적 관점에서 본 초복소량과 표현론.\\
Atti Congresso Bologna 2 (1928), 71--73쪽
\item[34.] 초복소량과 표현론.\\
Math. Z. 30 (1929), 641--692쪽
\item[35.] 정수적 함수로 이루어진 극대영역에 관하여.\\
Rec. Soc. Math. Moscou 36 (1929), 65--72쪽
\item[36.] 아이디얼 미분과 디퍼런트.\\
J. Ber. d. DMV 39 (1930), 17쪽
\item[37.] 고차 분기가 없는 체의 정규기저.\\
J. Reine Angew. Math. 167 (1932), 147--152쪽
\item[38.] R. 브라우어 및 H. 하세와 공동: 대수론의 한 주정리의 증명.\\
J. Reine Angew. Math. 167 (1932), 399--404쪽
\item[39.] 초복소계와 가환대수 및 수론의 관계.\\
Verhandl. Intern. Math.-Kongreß Zürich 1 (1932), 189--194쪽
\item[40.] 비가환 대수.\\
Math. Z. 37 (1933), 514--541쪽
\item[41.] 상대 갈루아 수체에 대한 주종정리.\\
Math. Ann. 108 (1933), 411--419쪽
\item[42.] 분할 교차곱과 그 극대 오더.\\
Actualités scientifiques et industrielles 148 (1934), 15쪽
\item[43.] 아이디얼 미분과 디퍼런트.\\
J. Reine Angew. Math. 188 (1950), 1--21쪽
\end{enumerate}

\noindent H. 카프페러, 뇌터의 대수함수 기본정리에 대한 필요충분 중복도 조건\\
(E. 뇌터와 공동으로 쓴 부록 포함).\\
Math. Ann. 97 (1927), 559--567쪽

\clearpage
\begin{center}
{\Large\bfseries 에미 뇌터의 단신 및 서평 목록}
\end{center}

\begin{center}
\emph{독일수학회 연보에 실린 단신}
\end{center}

\begingroup
\setlength{\parindent}{0pt}
\setlength{\parskip}{0.15em}
임의의 미분식의 불변식에 관하여, 27, 제II부 (1918), 28쪽\par
불변 변분문제, 27, 제II부 (1918), 47쪽\par
이원형식의 정수 계수 불변식의 유한성, 28, 제II부 (1919), 29쪽\par
정수 계수 다항식과 거듭제곱급수에 관하여, 28, 제II부 (1919), 29--30쪽\par
가군론의 분해정리, 29, 제II부 (1920), 54쪽\par
기본약수와 일반 아이디얼 이론, 30, 제II부 (1921), 32쪽\par
립시츠의 한 공식에 대한 지적, 30, 제II부 (1921), 48쪽\par
일반 아이디얼 이론에서 힐베르트 정리의 유사물, 32, 제II부 (1923), 20--21쪽\par
공리화 문제에 대한 군론적 논평, 33, 제II부 (1925), 21쪽\par
임의의 체에서의 갈루아 이론, 33, 제II부 (1925), 120쪽\par
아이디얼 이론의 힐베르트 수, 34, 제II부 (1926), 101쪽\par
아이디얼의 미분몫과 분기이론, 38, 제II부 (1929), 81쪽\par
\endgroup

\begin{center}
\emph{독일수학회 연보에 실린 서평}
\end{center}

\begingroup
\setlength{\parindent}{0pt}
\setlength{\parskip}{0.35em}
E. Study, \emph{Einleitung in die Theorie der Invarianten linearer Transformationen auf Grund der Vektorenrechnung}. 제1부. (Die Wissenschaft, Bd. 71) Braunschweig 1923, Vieweg.\\
36, 제II부 (1927), 71쪽\par
\emph{Leopold Kroneckers Werke}, 제4권. K. Hensel 편집. VI 및 509쪽. Leipzig 및 Berlin 1929, B. G. Teubner.\\
40, 제II부 (1931), 11--12쪽\par
\emph{Leopold Kroneckers Werke}, 제5권. K. Hensel 편집. 527쪽. Leipzig 및 Berlin 1930, B. G. Teubner.\\
41, 제II부 (1932), 17쪽\par
\emph{Leopold Kroneckers Werke}, 제III.2권. K. Hensel 편집. 215쪽. Leipzig 및 Berlin 1931, B. G. Teubner.\\
42, 제II부 (1933), 38--39쪽\par
D. E. Rutherford, \emph{Modular Invariants}. Cambridge Tracts No. 27. Cambridge 1932, VII 및 84쪽.\\
43, 제II부 (1934), 94--95쪽\par
\endgroup

\clearpage
\begin{center}
\emph{에미 뇌터가 편집에 참여한 책}
\end{center}

\begingroup
\setlength{\parindent}{0pt}
\setlength{\parskip}{0.5em}
R. Dedekind, \emph{Gesammelte mathematische Werke} I--III. R. Fricke, Emmy Noether 및 Øystein Ore 편집. Vieweg, Braunschweig 1930--1932\\
(제I권: 1930, 제II권: 1931, 제III권: 1932)\par
G. Cantor, \emph{Briefwechsel mit R. Dedekind}. Emmy Noether 및 J. Cavaillès 편집, 60쪽. Actualités scient. et ind. 581, Paris 1937\par
\endgroup
